\documentclass[12pt,a4paper]{article}

% --- PREAMBOLO E PACCHETTI ---
\usepackage[utf8]{inputenc}
\usepackage[italian]{babel}
\usepackage[T1]{fontenc}
\usepackage{amsmath, amssymb, amsthm}
\usepackage{graphicx}
\usepackage{xcolor}
\usepackage{tcolorbox}
\tcbuselibrary{skins, breakable, theorems}
\usepackage{listings}
\usepackage{geometry}
\usepackage{hyperref}
\usepackage{tikz}
\usetikzlibrary{shapes.geometric, arrows.meta, positioning, calc}

\geometry{margin=2cm}

% --- PALETTE COLORI ACCADEMICA ---
\definecolor{paletteIndigo}{HTML}{3F51B5}
\definecolor{paletteRegalia}{HTML}{522D80}
\definecolor{paletteLenurple}{HTML}{BA55D3}
\definecolor{paletteLavender}{HTML}{E6E6FA}

% --- CONFIGURAZIONE STILE CODICE (jsbox) ---
\lstdefinelanguage{JavaScript}{
  keywords={break, case, catch, class, const, continue, debugger, default, delete, do, else, enum, export, extends, false, finally, for, function, if, import, in, instanceof, new, null, return, super, switch, this, throw, true, try, typeof, var, void, while, with, yield, get, set, static, require},
  keywordstyle=\color{paletteIndigo}\bfseries,
  ndkeywords={boolean, number, String, Array, Date, Error, TypeError, ReferenceError, Map, Set, JSON, RegExp, Math, console, Buffer, process, fs, http, os},
  ndkeywordstyle=\color{paletteRegalia}\bfseries,
  identifierstyle=\color{black},
  sensitive=true,
  comment=[l]{//},
  morecomment=[s]{/*}{*/},
  commentstyle=\color{gray}\ttfamily,
  stringstyle=\color{paletteLenurple}\ttfamily,
  morestring=[b]',
  morestring=[b]"
}

\lstset{
  language=JavaScript,
  backgroundcolor=\color{paletteLavender!20},
  basicstyle=\small\ttfamily,
  numbers=left,
  numberstyle=\tiny\color{gray},
  stepnumber=1,
  numbersep=8pt,
  showstringspaces=false,
  breaklines=true,
  frame=leftline,
  bordercolor=\color{paletteIndigo},
  tabsize=2
}

% --- AMBIENTI BOX PERSONALIZZATI ---
\newtcolorbox{definition}[2][]{
  colback=paletteLavender!30, colframe=paletteIndigo, fonttitle=\bfseries,
  title=Definizione: #2, breakable, enhanced, attach boxed title to top left={xshift=3mm, yshift=-2mm},
  boxed title style={colback=paletteIndigo}
}

\newtcolorbox{concept}[2][]{
  colback=paletteLenurple!10, colframe=paletteLenurple, fonttitle=\bfseries,
  title=#2, breakable, enhanced
}

\newtcolorbox{exercise}[2][]{
  colback=white, colframe=paletteRegalia, fonttitle=\bfseries,
  title=Esercizio: #2, breakable, enhanced
}

\newtcolorbox{solution}[1][]{
  colback=green!5, colframe=green!40!black, fonttitle=\bfseries,
  title=Soluzione, breakable, enhanced
}

\title{\textcolor{paletteRegalia}{\Huge \textbf{Laboratorio I: La Bibbia Completa}} \\ 
\large \textit{Eccezioni, Classi di Sistema, Moduli ed Esercitazione GradeBook}}
\author{A.A. 2025/2026}
\date{}

\begin{document}

\maketitle
\tableofcontents
\newpage

\section{Bug, Errori ed Eccezioni}

Durante lo sviluppo di un programma possono essere introdotti dei difetti, comunemente definiti bug. 

\subsection{Tipologie di Errori}
Gli errori possono essere di diversa natura:
\begin{itemize}
    \item \textbf{Errori di sintassi}: Di solito riconosciuti subito dal sistema.
    \item \textbf{Errori di battitura (typo)}: Un nome di una variabile/funzione scritto in modo sbagliato possono generare un errore ma non sempre in JavaScript.
    \item \textbf{Errori di logica}: L'algoritmo proposto non funziona in alcuni casi; di solito sono silenziosi fino a quando non si incontra il caso limite.
    \item \textbf{Errori esterni al programma}: Errori di input/output (l'utente inserisce dei dati mal formattati), un file cancellato per sbaglio, fallimento del disco rigido, rete non funzionante per cavo rotto o router deceduto, ecc.
\end{itemize}

\subsection{Gestione degli Errori}
Il debugging per trovare i bug è un processo che può essere molto difficile. Scrivere codice ben organizzato e usare la modalità \texttt{'use strict'} ci aiuta a trovare i bug.
Alcuni tipi di errori non possono essere evitati. JavaScript prova sempre a trovare una soluzione, ma a volte restituisce un risultato sbagliato o un messaggio di errore. Non è una buona idea esibire i messaggi di errore all'utente. Bisogna gestire gli errori, dare feedback e continuare l'esecuzione.

\begin{definition}{Eccezione}
Meccanismo di interruzione del programma in caso di errore. Viene lanciata quando l'errore compare. Le eccezioni possono essere automatiche o lanciate manualmente al bisogno. Possono essere catturate nel posto più adatto del programma, anche solo per evitare messaggi di errore brutti.
\end{definition}

\subsection{Catturare un'Eccezione: try-catch-finally}
Il blocco \texttt{try} include i comandi che possono generare un errore. Il codice viene eseguito fino al primo comando che genera un errore (incluso nelle funzioni annidate). I comandi successivi all'errore non vengono eseguiti. Se compare un errore, si esegue il blocco \texttt{catch}. La variabile \texttt{e} contiene informazioni sull'errore. Se non c'è errore, il blocco \texttt{catch} viene saltato.

\begin{lstlisting}
try {
  // comandi
} catch (e) {
  // comandi gestione errore
} finally {
  // comandi finali
}
\end{lstlisting}

L'eccezione viene memorizzata in una variabile di tipo \texttt{Error} che espone proprietà come: \texttt{message}, \texttt{stack}, \texttt{name}. 
Il blocco \texttt{finally} permette di includere comandi da eseguire \textit{sempre}, sia nel caso di errore che non, utile per lasciare i dati in uno stato valido o chiudere file.

\subsection{Esempio Pratico: Creare e Lanciare Eccezioni Custom}
Possiamo definire un'intera gerarchia di errori per distinguere tra vari tipi di problematiche.

\begin{center}
\begin{tikzpicture}[
  node distance=1.5cm and 2cm,
  classnode/.style={rectangle, draw=paletteRegalia, fill=paletteLavender!40, thick, minimum width=3cm, minimum height=1cm, align=center, font=\bfseries}
]
\node[classnode] (error) {Error};
\node[classnode, below=of error] (expr) {ExpressionError};
\node[classnode, below left=of expr] (opnd) {OperandError};
\node[classnode, below right=of expr] (oprt) {OperatorError};

\draw[-Latex, thick, paletteIndigo] (expr) -- (error);
\draw[-Latex, thick, paletteIndigo] (opnd) -- (expr);
\draw[-Latex, thick, paletteIndigo] (oprt) -- (expr);
\end{tikzpicture}
\end{center}

Ecco il codice completo per l'implementazione della funzione \texttt{calc} e \texttt{calcAll} con la gerarchia di eccezioni:

\begin{lstlisting}
class ExpressionError extends Error { ; }
class OperatorError extends ExpressionError { ; }
class OperandError extends ExpressionError {
  operand;
  constructor(message, operand) {
    super(message);
    this.operand = operand;
  }
}

function calc(a) {
  var [op, ...x] = a;
  for (let n of x)
    if (typeof n != "number")
      throw new OperandError("Operand is not correct, should be number", n);

  switch (op) {
    case '+': return x.reduce((y,z)=>(y+z),0);
    case '-': 
      var [x1, ...x_rest] = x;
      return x_rest.reduce((y,z)=>(y-z),x1);
    case '*': return x.reduce((y,z)=>(y*z),1);
    case '/': 
      var [x1, ...x_rest] = x;
      return x_rest.reduce((y,z)=>(y/z), x1);
    default:
      throw new OperatorError("Operator is not correct, should be +*-/");
  }
}

function calcAll(expressions) {
  let res = [];
  for (let expr of expressions) {
    let val = undefined;
    try {
      val = calc(expr);
    } catch (e) {
      if (e instanceof ExpressionError) val = null;
      else throw e;
    }
    res.push(val);
  }
  return res;
}
\end{lstlisting}

\section{Classi di Sistema (Built-in) e Moduli Utili}

\subsection{Oggetti Predefiniti e Wrapper}
Lo standard JavaScript include oggetti predefiniti. Per ogni tipo base, esiste una classe wrapper con lo stesso nome con iniziale maiuscola: \texttt{Number}, \texttt{Boolean}, \texttt{String}, \texttt{BigInt}. I costruttori convertono l'argomento nel tipo corrispondente:
\begin{itemize}
    \item \texttt{Number("1.33") === 1.33}.
    \item \texttt{String(14) === "14"}.
    \item Offrono costanti e metodi statici, es. \texttt{Number.EPSILON}, \texttt{Number.isFinite(3.14)}.
\end{itemize}

\subsection{Date e RegExp}
\begin{concept}{Classe Date}
Fornisce oggetti che rappresentano date dal 20 Aprile 271.821 a.C. al 13 Settembre 275.760 d.C. 
Attenzione: \textbf{Il mese parte da 0 = Gennaio}.
Gli oggetti \texttt{Date} sono sostanzialmente numeri espressi in millisecondi dal 1 Gennaio 1970, consentendo operazioni aritmetiche (es. \texttt{b - a}).
\end{concept}

\begin{concept}{Espressioni Regolari (RegExp)}
Esprimono schemi di corrispondenza su stringhe. Hanno una denotazione letterale con \texttt{/.../}.
Metodi principali: \texttt{e.test(s)}, \texttt{e.exec(s)}, \texttt{s.match(e)}, \texttt{s.matchAll(e)}.
Flag: \texttt{/g} (tutte le occorrenze), \texttt{/i} (ignora maiuscole/minuscole), \texttt{/m} (multilinea).
\end{concept}

\subsection{Array di Macchina e Collezioni}
\begin{itemize}
    \item \textbf{Typed Array}: \texttt{Int8Array}, \texttt{Float64Array}, ecc. Manipolano blocchi di memoria di dimensione esatta, utili per scambiare dati in formati pre-determinati via rete o file.
    \item \textbf{Set}: Insiemi di valori qualunque. A differenza degli oggetti, le chiavi possono essere di qualunque tipo. Metodi: \texttt{add(e)}, \texttt{delete(e)}, \texttt{clear()}, \texttt{has(e)}, \texttt{size}.
    \item \textbf{Map}: Associa chiavi a valori. Metodi: \texttt{set(k,v)}, \texttt{delete(k)}, \texttt{get(k)}, \texttt{keys()}, \texttt{values()}.
\end{itemize}

\subsection{JSON}
Formato di serializzazione POJO (Plain Old JavaScript Object). Due metodi statici:
\begin{itemize}
    \item \texttt{JSON.stringify(o)}: da oggetto a stringa.
    \item \texttt{JSON.parse(s)}: da stringa a oggetto.
\end{itemize}

\subsection{Moduli Host-Defined in Node.js}
Node.js fornisce moduli built-in potentissimi per interagire con l'OS.

\subsubsection{Modulo \texttt{fs} (File System)}
Contiene funzioni per leggere/scrivere file e directory.

\begin{lstlisting}
const fs = require('fs');

// Lettura sincrona (adatta a file piccoli)
var s = fs.readFileSync('index.js');
console.log(s.toString());

// Lettura asincrona con callback (adatta a file grandi - chunk/stream)
fs.readFile('dante.txt', 'utf-8', (err, data) => {
  console.log(err, data.substr(0,100));
});
\end{lstlisting}

\subsubsection{Moduli \texttt{os}, \texttt{process}, \texttt{http}}
\begin{itemize}
    \item \textbf{\texttt{os}}: Info su macchina, OS, CPU (\texttt{os.cpus()}), memoria (\texttt{os.totalmem()}), architettura (\texttt{os.arch()}).
    \item \textbf{\texttt{process}}: Info sul processo in esecuzione, argomenti riga di comando, variabili d'ambiente (\texttt{process.env}).
    \item \textbf{\texttt{http}}: Per parsing URL e protocollo web, per creare server.
\end{itemize}

\begin{lstlisting}
const http = require('http');
var server = http.createServer(function (req, res) {
  res.writeHead(200, {'Content-Type': 'text/plain; charset=utf-8'});
  if (req.url.match(/dante/i)) {
    res.end(fs.readFileSync('dante.txt', 'utf-8'));
  } else {
    res.end("Circolare, non c'e niente da vedere!");
  }
});
server.listen(1337);
\end{lstlisting}

\section{Moduli in JavaScript}

I moduli servono per dividere il sorgente in più file per tenere in ordine, riutilizzare parti del programma e distribuire le responsabilità tra programmatori.

\begin{itemize}
    \item \textbf{CommonJS (Tradizionale)}: Uso di \texttt{require()} e \texttt{exports}.
\begin{lstlisting}
const mod = require('./modulo');
exports.fun = f2;
\end{lstlisting}
    \item \textbf{ES6 (Moderno)}: Uso di \texttt{import} ed \texttt{export}. Richiede la configurazione \texttt{"type": "module"} nel file \texttt{package.json}.
\begin{lstlisting}
import * as mod from './modulo.js';
export const PIGRECO = 4;
\end{lstlisting}
\end{itemize}

\section{Esercitazione Completa: Gestione Voti Universitari}

Di seguito tutti gli esercizi dal PDF LAB1\_29 e la soluzione integrale con classi ed eccezioni.

\begin{exercise}{1, 2, 3, 4, 5: Classi Grade, GradeBook e Errori}
Implementare il sistema di registrazione voti, definendo esami (\texttt{Grade}) con vincoli (voto tra 18 e 30, crediti positivi, controllo lode) e un registro (\texttt{GradeBook}) diviso in \texttt{BachelorGradeBook} (180 CFU) e \texttt{MasterGradeBook} (120 CFU). Prevedere eccezioni custom per errori di inserimento, crediti mancanti ed esami duplicati.
\end{exercise}

\begin{center}
\begin{tikzpicture}[
  node distance=1.5cm and 3cm,
  classnode/.style={rectangle, draw=paletteRegalia, fill=paletteLavender!40, thick, minimum width=4cm, minimum height=1cm, align=center, font=\bfseries}
]
\node[classnode] (gb) {GradeBook};
\node[classnode, below left=of gb] (bgb) {BachelorGradeBook};
\node[classnode, below right=of gb] (mgb) {MasterGradeBook};

\draw[-Latex, thick, paletteIndigo] (bgb) -- (gb);
\draw[-Latex, thick, paletteIndigo] (mgb) -- (gb);
\end{tikzpicture}
\end{center}

\begin{solution}[Codice Integrale Esercizi 1-5]
\begin{lstlisting}
class InvalidCreditsError extends Error {;}
class InvalidValueError extends Error {;}

class Grade {
  #name; #credits; #date; #grade; #honors;
  constructor (name, credits, date, grade, honors) {
    this.#name = name;
    if (credits <= 0) throw new InvalidCreditsError("Crediti negativi o uguali a 0");
    this.#credits = credits;
    this.#date = new Date(date);
    if (grade < 18 || grade > 30) throw new InvalidValueError("Voto fuori scala");
    this.#grade = grade;
    if (honors && grade != 30) throw new InvalidValueError("Lode non assegnabile");
    this.#honors = honors;
  }
  
  get name() { return this.#name; }
  get credits() { return this.#credits; }
  get date() { return this.#date; }
  get grade() { return this.#grade; }
  get honors() { return this.#honors; }

  equals(g) {
    return this.name == g.name && this.credits == g.credits && 
           this.date - g.date == 0 && this.grade == g.grade && this.honors == g.honors;
  }

  toString() {
    let s = this.name + "(" + this.credits + "CFU):" + this.grade;
    if (this.honors) s += "L";
    return s;
  }
}

class DuplicatedGradeError extends Error {; }
class MissingCreditsError extends Error {; }

class GradeBook {
  #grades; #totalCredits;
  
  constructor (totalCredits) {
    this.#grades = [];
    this.#totalCredits = totalCredits;
  }
  
  get totalCredits() { return this.#totalCredits; }
  
  get credits() {
    let c = 0;
    for (let g of this.#grades) c += g.credits;
    return c;
  }
  
  get missingCredits() { return this.#totalCredits - this.credits; }
  
  get average() {
    let sum = 0; let count = 0;
    for (let g of this.#grades) {
      sum += (g.honors ? 32 : g.grade) * g.credits;
      count += g.credits;
    }
    return sum / count;
  }
  
  get startingGrade() {
    if (this.missingCredits != 0) throw new MissingCreditsError("Crediti mancanti");
    return Math.round(this.average * 11 / 3);
  }

  #search(grade) {
    for (let g of this.#grades)
      if (g.equals(grade)) return true;
    return false;
  }

  register(grade) {
    if (!(grade instanceof Grade)) throw new TypeError("Non e un voto");
    if (this.#search(grade)) throw new DuplicatedGradeError("Esame gia registrato");
    this.#grades.push(grade);
  }

  toString() {
    let s = "VOTI:\n";
    for (let g of this.#grades) s += "\t" + g.toString() + "\n";
    return s;
  }
}

class BachelorGradeBook extends GradeBook { constructor() { super(180); } }
class MasterGradeBook extends GradeBook { constructor() { super(120); } }

// Test:
let gb = new GradeBook(180);
gb.register(new Grade("LAB 1", 12, "2026/02/26", 30, true));
\end{lstlisting}
\end{solution}

\begin{exercise}{Esercizi 6-9: Persistenza, Moduli e Regex (Per Casa)}
\begin{itemize}
    \item \textbf{Esercizio 6}: Aggiungere \texttt{exportJSON(f)} e \texttt{fromJSON(f)} usando File System \texttt{fs}.
    \item \textbf{Esercizio 7}: Creare modulo \texttt{gradebook.js} che esporti le classi.
    \item \textbf{Esercizio 8}: Validazione codice matricola (esattamente 6 cifre) usando le RegExp.
    \item \textbf{Esercizio 9}: Validazione email studenti (\texttt{n.cognome[n]@studenti.unipi.it}).
\end{itemize}
\end{exercise}

\begin{solution}[Codice Esercizi 8 e 9]
\begin{lstlisting}
// Esercizio 8
function matricolaValida(m) {
  return /^\d{6}$/.test(m);
}

// Esercizio 9
function emailStudenteUnipiValida(e) {
  return /^[a-z]\.[a-z]+[1-9]?\d*@studenti\.unipi\.it$/.test(e);
}
\end{lstlisting}
\end{solution}

\end{document}
